\scanpage{24}

\[
G_{A'}\simeq\mathcal H_{A'}
\]
Comme invariant ($C$ \unsure{donc} $G_{A'}$ \inserted{de $G_{A'}$} étant ici fixée) on trouve que cette donnée \inserted{à isom près} représente un point d'un \uuline{torseur} sur
\[
H^1(A',\mu_2)\simeq\prod_{\ell\ne p,\infty}\underbrace{\mathbb Z_\ell^*/\mathbb Z_\ell^{*2}}_{\begin{gathered}\simeq\mathbb Z/2\mathbb Z\text{ si }\ell\ne2\\\simeq\mathbb Z/2\mathbb Z\times\mathbb Z/2\mathbb Z\text{ si }\ell=2\end{gathered}}
\]
\deleted{\[
\unsure{H^2(\mathbb Q\to A',\mu_2)\simeq\operatorname{Coker}(}
\]}
\margin{\uline{NB} il n'y a absolument pas de façon naturelle de choisir une origine \inserted{de façon ``\unsure{abstraite}''} \deleted{\unsure{voir dans cohomologie} \unsure{par des classes} \unsure{dans} $\unsure{\mu_\ell}$ mod $(\mathbb Z/\ell\mathbb Z)^{*2}$ $(\ell\ne2)$ $\unsure{\mu_8}$ mod $(\mathbb Z/8\mathbb Z)^{*2}$}}

Si on tient compte du fait que la gerbe $G_{A'}$ est déterminée seulement à équivalence (non unique près), le choix d'une \unsure{classe d'isomorphie d'équivalence permet} \unsure{diffère} donc d'une classe de
\[
H^1(\underbrace{\mathbb Z[\tfrac1p]}_{A'},\mu_2)\simeq\mathbb Z/2\mathbb Z\times\mathbb Z/2\mathbb Z
\]
(\unsure{l'élt} \deleted{un él de $\mathbb Z/2\mathbb Z$} \inserted{$(1,0)$ resp. $(0,1)$} correspond au rev. de groupe $\mu_2$ obtenu par $\sqrt[2]{p}$ \inserted{resp. $\sqrt[2]{-1}$}) on trouve un invariant \deleted{dans} dans un quotient
\[
\Pi/\rho(\mathbb Z/2\mathbb Z),
\]
où $\Pi$ est le torseur sur $H^1(A',\mu_2)=\prod_{\ell\ne p,\infty}\mathbb Z_\ell^*/\mathbb Z_\ell^{*2}$ des trivialisations de $\mathcal H_{A'}$, et
\[
\rho:H^1(\mathbb Z[\tfrac1p],\mu_2)\simeq(\mathbb Z/2\mathbb Z)\inserted{\times\mathbb Z/2\mathbb Z}\longrightarrow H^1(A',\mu_2)
\]
est l'hom naturel, défini par la condition
\[
\rho(1,0)_\ell\inserted{\text{(resp. }\rho(0,1)_\ell\text{)}}=
\begin{cases}
0&\text{si }p\inserted{\text{(resp. }-1\text{)}}\text{ un carré dans }\mathbb Z_\ell^*\quad(\overset{\inserted{\text{si }\ell\ne2}}{\Longleftrightarrow}p\text{ carré dans }(\mathbb Z/\ell\mathbb Z)^*\text{) si }\deleted{\ell\ne2)}\\
1&\text{sinon}
\end{cases}
\]
\margin{on peut éliminer ce groupe $\mathbb Z/2\mathbb Z\times\mathbb Z/2\mathbb Z$ grâce à la donnée de la structure \unsure{suppl.} \unsure{a)} polarisation à l'$\infty$ b) ``cristalline'' en $p$}
\figure{\begin{tikzpicture}[x=1cm,y=1cm]\draw (0.04,0.34) -- (0.08,0.17);\draw (0,0.20) .. controls (-0.03,0.10) and (0.06,0.03) .. (0.10,0);\end{tikzpicture}}

\mspage{23}
